\documentclass[11pt]{article}
\usepackage[margin=0.7in]{geometry}
\usepackage[T1]{fontenc}
\usepackage{lmodern}
\usepackage{parskip}
\usepackage{enumitem}
\setlist[itemize]{topsep=6pt,itemsep=6pt,parsep=2pt,partopsep=0pt,leftmargin=1.2em}
\pagestyle{plain}

\title{DGGR Final Presentation Notes}
\date{}

\begin{document}
\maketitle

Simple glanceable notes for the 10-slide main deck.

\section*{Slide 1: Deep Generative Genre Remastering}
\begin{itemize}
\item Open by contrasting coat-of-paint style transfer with controlled genre remastering.
\item Frame the project as: separate content and style, build a target space, generate toward that target, then make it survive long-form playback.
\item Point to the pipeline strip and the content / realism / style tradeoff.
\item Keep this slide conceptual and do not play audio yet.
\item End by saying the first two labs matter because without trustworthy control, later audio is hard to interpret.
\end{itemize}

\section*{Slide 2: Lab 1}
\begin{itemize}
\item Explain that Lab 1 made the repo technically trustworthy by learning \texttt{z\_content}, \texttt{z\_style}, and a music gate.
\item Phrase it simply: content is what the song is doing, style is how it sounds.
\item Walk through the encoder split diagram while explaining the three outputs.
\item Style probe 0.9417 means style is really concentrated in the style branch.
\item Leakage 0.1083 means content is not secretly carrying style under another name.
\item Gate ROC-AUC 0.9299 means the model learned musically valid structure well enough to trust later stages.
\item End by saying that once style existed as a real signal, the next question was whether it could become a stable target.
\end{itemize}

\section*{Slide 3: Lab 2}
\begin{itemize}
\item Explain that Lab 2 turned style from an embedding into a target space the generator could actually aim for.
\item Point to the cluster figure first, then the centroid or target-vector diagram.
\item Silhouette 0.4939 means the space is not one tangled cloud.
\item Linear probe 0.8554 means the genre signal is accessible instead of buried.
\item Nearest-centroid 0.8514 means the centroids behave like real genre anchors.
\item End by saying that once representation and target space were in place, the remaining question became generation.
\end{itemize}

\section*{Slide 4: From Codec Editing to Diffusion Generation}
\begin{itemize}
\item Introduce codec as the first strong proof that controlled reconstruction and editing were possible.
\item Then explain that diffusion became the practical backbone because it gave the project a stronger realism anchor.
\item Point to the side-by-side codec versus diffusion comparison.
\item Play one short codec clip and one short diffusion clip.
\item End by saying that once diffusion became the main path, the core tradeoff became obvious.
\end{itemize}

\section*{Slide 5: Short-Form Generation}
\begin{itemize}
\item Use this slide to make the content-versus-style tradeoff feel obvious.
\item Say clearly that style under-shoot is disappointing, but identity collapse is destructive.
\item Point to the conservative / balanced / over-pushed comparison.
\item Play a safe but style-light clip, then a balanced clip, then an over-pushed clip.
\item End by saying that even promising local clips still had to survive time.
\end{itemize}

\section*{Slide 6: Long-Form Generation}
\begin{itemize}
\item Explain that long-form accompaniment generation is the real remaining bottleneck.
\item Say that local quality is not enough because seams, warble, drift, and instability build up over time.
\item Point to the chunk timeline or seam diagram.
\item Play one local-good clip and then the matching long-form failure.
\item End by saying that this is why the best practical system came from engineering around the generator.
\end{itemize}

\section*{Slide 7: Hybrid Engineering}
\begin{itemize}
\item Present the hybrid path as a systems solution, not a compromise in the bad sense.
\item Explain that we preserved vocals where the model was fragile and generated accompaniment where style mattered most.
\item Point to the preserved-vocal hybrid flow.
\item Play the before / after practical comparison and then the best practical clip.
\item End by saying that once we had a practical winner, the next question was whether later branches could beat it.
\end{itemize}

\section*{Slide 8: Later Experiments}
\begin{itemize}
\item Frame this as the learning slide, not the graveyard slide.
\item Explain that stronger style movement and new local textures were possible, but long-form stability stayed hard.
\item Point to the compact family table and summarize what each knob improved or broke.
\item Play only one or two representative later-branch clips.
\item End by saying the conclusion is not that nothing worked, but that the hard part became clearer.
\end{itemize}

\section*{Slide 9: What We Actually Solved}
\begin{itemize}
\item Keep this slide short and confident.
\item Say that the project succeeded as a pipeline even though it did not fully solve long-form generation.
\item State the four defendable wins: content/style representation, target-space control, locally plausible accompaniment generation, and a usable hybrid workflow.
\item End by saying that this leaves one final question: what exactly still remains?
\end{itemize}

\section*{Slide 10: What Still Remains}
\begin{itemize}
\item End on the real unresolved problem: long-form accompaniment generation.
\item Say that realism, style separation, new instrumental character, and song identity still need to survive together over time.
\item Do not play audio here; let the conclusion land cleanly.
\item Final line: the next push should target long-form stability before even harder style pressure.
\end{itemize}

\end{document}
